% fixture: display/align-amsmath
% tier: T3
% engine: document
% covers: amsmath align and align*: \measure@ (column widths, \totwidth@), \eqnshift@ and \alignsep@ for one and two column pairs, \minalignsep, tags placed by \place@tag, \nonumber, empty fields, \openup\jot and \interdisplaylinepenalty, inside a paragraph and opening one
% issue: #30
\documentclass{article}
\usepackage{amsmath}
\begin{document}
Aligned equations are set by an alignment inside a display: right-aligned and left-aligned fields alternate, the block is centred, and each row's number sits at the right margin.
\begin{align}
a &= b + c \\
d &= e + f \nonumber \\
\sum_{i=1}^n i &= \frac{n(n+1)}{2}
\end{align}
and two pairs of columns:
\begin{align}
x &= 1 & y &= 2 \\
x^2 + y^2 &= 5 & & = z
\end{align}
Aligned equations are set by an alignment inside a display: right-aligned and left-aligned fields alternate, the block is centred, and each row's number sits at the right margin.

\begin{align*}
f(x) &= x^2 \\
&= x \cdot x
\end{align*}
Aligned equations are set by an alignment inside a display: right-aligned and left-aligned fields alternate, the block is centred, and each row's number sits at the right margin.
\end{document}
